\documentclass[border=0pt]{standalone}
% The shared preamble of every docs diagram (docs/diagrams/<name>.tex).
% dvisvgm's driver first, so patterns and opacity survive into the SVG.
\def\pgfsysdriver{pgfsys-dvisvgm.def}
\usepackage[siunitx,RPvoltages]{circuitikz}
\usepackage{xcolor}
% Latin Modern: outline fonts for every glyph a label uses, µ included.
\usepackage[T1]{fontenc}
\usepackage{lmodern}
\usetikzlibrary{arrows.meta,backgrounds,calc,patterns,positioning}
\tikzset{
  % A white card, so diagrams read in mdBook's dark themes too.
  every picture/.append style={show background rectangle,
    background rectangle/.style={fill=white,inner sep=4pt}},
  % Layout geometry: layers told apart by fill and hatching, never colour alone.
  layer a/.style={draw=black,line width=0.6pt,fill=gray!20},
  layer b/.style={draw=black,line width=0.6pt,fill=gray!45},
  layer c/.style={draw=black,line width=0.6pt,fill=white,postaction={pattern=north east lines}},
  layer d/.style={draw=black,line width=0.6pt,fill=white,postaction={pattern=crosshatch dots}},
  marker/.style={draw=black,dashed,line width=0.6pt},
  % A measured distance and its label.
  dim/.style={{Latex[length=2mm]}-{Latex[length=2mm]},line width=0.5pt},
  dim label/.style={font=\small,fill=white,inner sep=1pt},
  % Where a rule fails: thick, dashed and labelled, not only red.
  violation/.style={draw=red!75!black,line width=1.4pt,dash pattern=on 3pt off 2pt},
  violation label/.style={font=\small\bfseries,text=red!75!black},
  caption/.style={font=\small},
}

\begin{document}
\begin{circuitikz}[x=0.7cm,y=0.7cm]
  % 1 unit = 1 um. A ring (0,0)-(3,2) with a hole (1,0.6)-(2,1.4), and an
  % L shape (4,0)-(7,0.8) plus (4,0.8)-(5,2).
  \def\ring{(0,0) rectangle (3,2) (1,0.6) rectangle (2,1.4)}
  \def\ell{(4,0) -- (7,0) -- (7,0.8) -- (5,0.8) -- (5,2) -- (4,2) -- cycle}
  \begin{scope}
    \fill[layer a,even odd rule] \ring; \fill[layer a] \ell;
    \node[caption,anchor=north] at (3.5,-0.5) {drawn: met1};
  \end{scope}
  \begin{scope}[xshift=6.3cm]
    \draw[marker] \ring; \draw[marker] \ell;
    \fill[layer b] (1,0.6) rectangle (2,1.4);
    \node[caption,anchor=north] at (3.5,-0.5) {\texttt{met1.holes()}};
  \end{scope}
  \begin{scope}[xshift=12.6cm]
    \fill[layer b] (0,0) rectangle (3,2); \fill[layer b] (4,0) rectangle (7,2);
    \draw[marker] (1,0.6) rectangle (2,1.4); \draw[marker] (5,0.8) -- (7,0.8);
    \node[caption,anchor=north] at (3.5,-0.5) {\texttt{met1.extents()}};
  \end{scope}
  \begin{scope}[xshift=18.9cm]
    \draw[marker] \ring; \draw[marker] \ell;
    \fill[layer b] (0,-0.2) rectangle (3,0); \fill[layer b] (0,2) rectangle (3,2.2);
    \fill[layer b] (4,-0.2) rectangle (7,0);
    \draw[line width=1.4pt] (0,0) -- (3,0) (0,2) -- (3,2) (4,0) -- (7,0);
    \node[caption,anchor=north,align=center] at (3.5,-0.5) {\texttt{met1.edges().horizontal()}\\\texttt{.with\_length(>= 3um)}\\\texttt{.extended(\dots, out: 200nm, in: 0nm)}};
  \end{scope}
\end{circuitikz}
\end{document}
